\section{The exact third-row common-group obstruction}
\label{sec:n3-character}

Throughout this section the common geometric source group is taken over
\(K\):
\[
G_3=\Dih(C_9),\qquad |G_3|=18.
\]
It has one identity, nine reflections of order two, two nonidentity
rotations of order three, and six rotations of order nine.  Let
\(\varepsilon\) be the reflection-sign representation.  For
\(1\le j\le4\), let \(U_j\) denote the two-dimensional irreducible with
\begin{equation}
\chi_{U_j}(r^k)=\zeta_9^{jk}+\zeta_9^{-jk},\qquad
\chi_{U_j}(r^ks)=0.
\label{eq:dihedral-irreducibles}
\end{equation}

\subsection{The Cayley quotient and residue action}

Nagel's complete-intersection identification is
\[
H_{\mathrm{var}}^{n-p,p}(X)\cong R_{p,d(X)},\qquad
d(X)=\sum_i d_i-n-r-2
\]
\cite[Proposition 2.16]{Nagel1997}.  In the present \((2,3)\)
threefold, \(n=3\), \(r=1\), and \(\sum_i d_i=5\), so
\(d(X)=-1\).  With \(p=1\), this gives the precise specialization
\begin{equation}
H^{2,1}(X_3)=H_{\mathrm{var}}^{2,1}(X_3)
\cong R_{1,-1}.
\label{eq:nagel-specialization}
\end{equation}
The \(\PP^7\) fivefold calculation in
\cite[\S5.4]{FaveroIlievKatzarkov2014} is useful context for the
Cayley bigrading, but it is not the locator for
\eqref{eq:nagel-specialization}.

For the Cayley polynomial \(yC_3+zQ_{3,\rho}\), bidegree
\((1,-1)\) has the 27 monomials
\[
zx_i\quad(0\le i<6),\qquad
yx_ix_j\quad(0\le i\le j<6).
\]
The relevant Jacobian subspace is generated by the six relations
\begin{equation}
3yx_i^2+z\frac{\partial Q_{3,\rho}}{\partial x_i},
\qquad 0\le i<6,
\label{eq:cayley-derivatives}
\end{equation}
and by \(yQ_{3,\rho}\).  These seven vectors are linearly independent over
\(K\), so
\begin{equation}
\dim R_{1,-1}=27-7=20.
\label{eq:cayley-dimension}
\end{equation}

If \(M_g\) is the coordinate matrix and \(A_g\) is the induced matrix on
the two-dimensional equation space, polynomial pullback alone is not the
cohomological action.  The residue action is
\begin{equation}
T_g=\frac{\det M_g}{\det A_g}\,T_g^{\mathrm{poly}}.
\label{eq:residue-correction}
\end{equation}
The determinant ratio is essential: omitting or inverting it changes the
character.

\begin{theorem}[Exact third-row character]
\label{thm:n3-character}
The residue-corrected action on \(H^{2,1}(X_3)\) has rotation traces
\[
\bigl(\Tr(r^k)\bigr)_{k=0}^{8}
=(20,-1,-1,2,-1,-1,2,-1,-1)
\]
and trace \(-2\) on every reflection.  Consequently
\begin{align}
H^{2,1}(X_3)
&=2\varepsilon+2U_1+2U_2+3U_3+2U_4,
\label{eq:h21-decomposition}\\
\Opack_3
&=4\varepsilon+4U_1+4U_2+6U_3+4U_4.
\label{eq:odd-decomposition}
\end{align}
For the Fermat rail, including its additional trivial line,
\[
\bigl(\Tr(r^k\mid\Epack_3)\bigr)_{k=0}^{8}
=(23,2,2,-4,2,2,-4,2,2),
\]
every reflection has trace \(-1\), and
\begin{equation}
\Epack_3
=\one+2\varepsilon+3U_1+3U_2+U_3+3U_4.
\label{eq:even-decomposition}
\end{equation}
\end{theorem}

\begin{proof}
The matrices induced by \(r\) and \(s\) preserve the seven-dimensional
relation space generated by \eqref{eq:cayley-derivatives} and
\(yQ_{3,\rho}\).  Taking the trace on the 27-dimensional monomial space,
subtracting the trace on that relation space, and multiplying by the residue
factor \eqref{eq:residue-correction} gives the displayed rotation and
reflection traces.  Inner products with the six irreducible characters
\(\one,\varepsilon,U_1,\ldots,U_4\) give
\eqref{eq:h21-decomposition}.  The character is real, and adjoining its
Hodge conjugate doubles the multiplicities, yielding
\eqref{eq:odd-decomposition}.

For the Fermat rail, the primitive Jacobian ring is
\[
K[x_0,\ldots,x_5]/(x_0^2,\ldots,x_5^2).
\]
The relevant squarefree degrees are \(0,3,6\).  Fixed-monomial traces,
multiplied by the residue factor \(\det M_g\), and the extra trivial line
give the second trace vector.  The same character inner products yield
\eqref{eq:even-decomposition}.  The exact trace-to-multiplicity ledger is
recorded in Appendix~\ref{app:proof-ledgers}.
\end{proof}

\begin{corollary}[No common central sector clears \(4/3\)]
\label{cor:no-central-sector}
No nonzero central \(G_3\)-isotypic summand of
\(\Epack_3\oplus\Opack_3\) has multiplicity divisible by three on both
pure rails.
\end{corollary}

\begin{proof}
The multiplicities are
\begin{center}
\begin{tabular}{@{}c|rrrrrr@{}}
\toprule
sector & \(\one\) & \(\varepsilon\) & \(U_1\) & \(U_2\) & \(U_3\) & \(U_4\)\\
\midrule
\(\Epack_3\) & \(1\) & \(2\) & \(3\) & \(3\) & \(1\) & \(3\)\\
\(\Opack_3\) & \(0\) & \(4\) & \(4\) & \(4\) & \(6\) & \(4\)\\
\bottomrule
\end{tabular}
\end{center}
Scaling an isotypic summand by \(4/3\) is integral only if its
multiplicity is divisible by three.  Every column fails this test on at
least one rail, so no nonempty collection of central isotypic sectors works.
Over the rational coefficient field, \(U_1,U_2,U_4\) form one Galois
orbit block.  Its multiplicity pair remains \((3,4)\), and therefore
coefficient-field packaging does not remove the obstruction.
\end{proof}

\begin{remark}[Common group versus common rational form]
\label{rem:common-group}
Theorem~\ref{thm:n3-character} is first a theorem for the common geometric
\(G_3\)-action over \(K\).  The Fermat packet in its standard rational
form and the complete intersection with \(M_3\)-descent need not carry the
same rational group scheme.  A formulation under one
\(\Gscheme_3\) uses the \(M_3\)-twisted rational form of the Fermat rail.
That twist preserves the split trace and the \(K\)-character above, while
its inert traces may differ.
\end{remark}
